% SOURCE_AUTHORITY: sga5_fr_workpass.tex, lines 13049--13084 (LF-normalized unit).
% SOURCE_SHA256: 791F4EFFC5E02832D5D77ED03518C8156D6F07E4C8238B03545DB93D883FBB28
\[
\begin{aligned}
\clK{\Delta}(\RG_C(F))
&=\operatorname{hom}\Bigl(
\clKstar{\Lambda[G]}\bigl(\RG_{C'}(\Lambda_{U',C'})\bigr),
\cl_{K(D(A[G])_\Delta)}(F_{\bar\eta})\Bigr)
 +\sum_{y\in Y}\clK{\Delta}(F_y)\\
&=\sum_{y\in Y}\clK{\Delta}(F_y)\\
&\quad +\operatorname{hom}\Bigl(
\EP(C)\clKstar{\Lambda[G]}(\Lambda[G])
 -\sum_{y\in Y}\bigl(Sw_y^{\Lambda}+\clKstar{\Lambda[G]}(\Lambda[G])\bigr),\\
&\hspace{10em}\cl_{K(D(A[G])_\Delta)}(F_{\bar\eta})\Bigr)\\
&=\sum_{y\in Y}\clK{\Delta}(F_y)
 +\EP(C)\clK{\Delta}F_{\bar\eta}
 -\sum_{y\in Y}\bigl(\alpha_y^\Delta(F)+\clK{\Delta}(F_{\bar\eta})\bigr)\\
&=\EP(C)\clK{\Delta}F_{\bar\eta}
 -\sum_{y\in Y}\varepsilon_y^\Delta(F).
\end{aligned}
\]
Tendo em conta que $\varepsilon_x^\Delta(F)=0$ para $x\in\mathcal U$, esta não é senão a fórmula (7.2). Para um complexo de feixes $F$ que satisfaça as condições de 7.1, introduziremos a notação
\[
\EP_\Delta(C,F)=\clK{\Delta}(\RG_C(F)).
\]

\begin{statement}{Corolário 7.9}
Se $F$, $F'$ são complexos de feixes que satisfazem as condições de 7.1 e se são localmente isomorfos (para a topologia étale) como complexos de $A$-módulos, tem-se:
\[
\EP_\Delta(C,F)=\EP_\Delta(C,F').
\tag{7.10}
\]
Em particular, se $F$ é localmente constante de valor $M^\bullet$, tem-se:
\[
\EP_\Delta(C,F)=\EP(C)\clK{\Delta}(M^\bullet).
\tag{7.11}
\]
\end{statement}
